\section{The normalized genus-one determinant}
\label{v4-arith:sc:determinant}

Let $Z$ contain each zero $\rho$ of $\xi$ with its multiplicity.
On $K_Z=\ell^2(Z)$, define
\[
 (\Theta_Zu)_\rho=\rho u_\rho,\qquad
 \operatorname{Dom}\Theta_Z=\{u:\sum_\rho|\rho u_\rho|^2<\infty\}.
\]
This is a closed normal operator. Its inverse is bounded, compact,
and Hilbert--Schmidt: $\xi(0)=1/2$, the zeros are discrete, and
$\sum_\rho|\rho|^{-2}<\infty$.

\begin{theorem}[complete determinant normalization]
\label{v4-arith:sc:normalized-det}
Define the second regularized Fredholm determinant of the diagonal
Hilbert--Schmidt operator $-s\Theta_Z^{-1}$ by
\[
 D_2(s)=\det{}_2(1-s\Theta_Z^{-1})
       :=\prod_{\rho\in Z}(1-s/\rho)e^{s/\rho}.
\]
With $B=-1-\gamma/2+\tfrac12\log(4\pi)$, set
\begin{equation}
 \Delta_Z(s):=\tfrac12e^{Bs}D_2(s).
 \label{v4-arith:sc:normalized-det-formula}
\end{equation}
This entire function equals $\xi(s)$.
It is the unique entire function of order at most one with this
zero divisor, $\Delta_Z(0)=1/2$, and $\Delta_Z'(0)/\Delta_Z(0)=B$.
Off the zero divisor it satisfies
\[
 (\log\Delta_Z)''(s)
       =-\operatorname{Tr}(\Theta_Z-s)^{-2}.
\]
The trace is an ordinary absolutely convergent trace.
Unitary conjugacy of $\Theta_Z$, including its domain, preserves
this determinant prescription.
\end{theorem}
\begin{proof}
For $|w|\le1/2$, $\log(1-w)+w=O(|w|^2)$.
The inverse-square sum proves normal convergence of the product
on every compact set, independently of the enumeration of zeros.
Its value and derivative at zero are $1$ and $0$.
Proposition~\ref{v4-arith:pm:constant} gives the product for $\xi$
with exactly the factor $e^{Bs}/2$. This proves the identity.
Hadamard factorization writes every other entire function of order
at most one with this divisor as $e^{a+bs}D_2(s)$.
The two initial conditions give $e^a=1/2$ and $b=B$.
Twice differentiating the normally convergent logarithmic product
gives $-\sum_\rho(\rho-s)^{-2}$.
The inverse-square estimate proves absolute trace convergence.
Unitary conjugacy preserves the diagonal eigenvalues and their
multiplicities, hence every factor in the prescribed product.
\end{proof}

\begin{proposition}[the divisor leaves an exponential ambiguity]
\label{v4-arith:sc:det-ambiguity}
For every $c\in\mathbb C$, the function $e^{cs}\xi(s)$ has the same
divisor, the same value at zero, and order at most one.
Its logarithmic derivative at zero is $B+c$.
Thus a divisor condition and a value normalization do not imply
the determinant identity. Both normalizations in
\eqref{v4-arith:sc:normalized-det-formula} are necessary.
\end{proposition}
\begin{proof}
The exponential has no zero, has value one at zero, and contributes
$c$ to the logarithmic derivative. Multiplication preserves the
stated order bound.
\end{proof}

The prescription above uses no complex powers of $s-\Theta_Z$.
A determinant defined by such powers requires a spectral cut and
meromorphic continuation of its trace. Agreement with $\Delta_Z$
requires the same second logarithmic derivative and the two stated
initial conditions. These conditions determine equality locally
near zero by integration, and globally by analytic continuation.
